\section{Complete Database Schema}
\label{app:ddl}

This appendix gives the version-1 DDL in full. Listings~\ref{lst:ddl-asset}
through \ref{lst:ddl-rend} in Section~\ref{sec:schema} reproduce the four
central tables with their design rationale; they are repeated here only where
necessary for the schema to be complete and executable as written. The text
below is the literal content of \texttt{ddlV1} referenced by the migration
registration in Section~\ref{sec:schema}.

\subsection{Connection Configuration}

\begin{lstlisting}[language=SQL, caption={Pragmas applied on every connection.}, label={lst:ddl-pragma}]
PRAGMA journal_mode  = WAL;         -- concurrent reads during autosave
PRAGMA synchronous   = NORMAL;      -- WAL makes FULL unnecessary
PRAGMA foreign_keys  = ON;          -- off by default in SQLite; must be set
PRAGMA cache_size    = -20000;      -- 20 MB, negative unit is kibibytes
PRAGMA busy_timeout  = 5000;
PRAGMA temp_store    = MEMORY;
PRAGMA mmap_size     = 268435456;   -- 256 MB
PRAGMA auto_vacuum   = INCREMENTAL; -- rendition churn would otherwise bloat
\end{lstlisting}

\texttt{foreign\_keys} deserves emphasis: SQLite disables foreign key
enforcement by default, per connection, and a schema full of \texttt{REFERENCES}
clauses that are never enforced is a schema that lies. GRDB applies this in the
connection preparation block so that no code path can open an unconfigured
connection.

\subsection{Metadata and Settings}

\begin{lstlisting}[language=SQL, caption={Schema metadata and key/value settings.}, label={lst:ddl-meta}]
-- Written by the migrator; read by diagnostics and by the crash reporter so
-- that a bug report identifies the exact schema and recipe versions.
CREATE TABLE schema_meta (
    key             TEXT PRIMARY KEY,
    value           TEXT NOT NULL
) STRICT;

INSERT INTO schema_meta(key, value) VALUES
    ('schema_version', '1'),
    ('recipe_version', '4'),
    ('created_by',     '1.0.0');

-- Application settings that are not user documents. Deliberately not in
-- UserDefaults: these participate in the same backup and migration story
-- as the library, and a setting that survives a restore while its library
-- does not is a support incident.
CREATE TABLE setting (
    key             TEXT PRIMARY KEY,
    value           TEXT NOT NULL,
    updated_at      REAL NOT NULL
) STRICT;

-- Per-device calibration constant g_cal relating the decoded middle-gray
-- value to the working-space unit. Device-specific, so it is keyed by
-- model identifier and never synchronized between devices.
CREATE TABLE device_calibration (
    model_id        TEXT NOT NULL,        -- e.g. "iPhone17,2"
    lens_id         TEXT NOT NULL,        -- wide | ultrawide | tele
    g_cal           REAL NOT NULL,
    method          TEXT NOT NULL,        -- 'shipped' | 'measured'
    measured_at     REAL,
    PRIMARY KEY (model_id, lens_id)
) STRICT;
\end{lstlisting}

\subsection{Library Tables}

\begin{lstlisting}[language=SQL, caption={Roll, asset, and tagging.}, label={lst:ddl-library}]
CREATE TABLE roll (
    id                TEXT PRIMARY KEY,
    name              TEXT NOT NULL,
    created_at        REAL NOT NULL,
    default_recipe_id TEXT REFERENCES recipe(id) ON DELETE SET NULL,
    notes             TEXT
) STRICT;

CREATE TABLE asset (
    id              TEXT PRIMARY KEY,
    roll_id         TEXT REFERENCES roll(id) ON DELETE SET NULL,
    file_relpath    TEXT,
    ph_local_id     TEXT,
    content_hash    TEXT NOT NULL,
    original_uti    TEXT NOT NULL,
    pixel_width     INTEGER NOT NULL,
    pixel_height    INTEGER NOT NULL,
    captured_at     REAL,
    imported_at     REAL NOT NULL,
    camera_model    TEXT,
    lens_model      TEXT,
    iso             INTEGER,
    exposure_time   REAL,
    f_number        REAL,
    focal_length    REAL,
    gps_lat         REAL,
    gps_lon         REAL,
    is_favorite     INTEGER NOT NULL DEFAULT 0,
    deleted_at      REAL,
    CHECK ((file_relpath IS NULL) <> (ph_local_id IS NULL))
) STRICT;

CREATE INDEX idx_asset_roll        ON asset(roll_id, captured_at DESC);
CREATE INDEX idx_asset_captured    ON asset(captured_at DESC)
                                   WHERE deleted_at IS NULL;
CREATE UNIQUE INDEX idx_asset_hash ON asset(content_hash)
                                   WHERE deleted_at IS NULL;
CREATE INDEX idx_asset_deleted     ON asset(deleted_at)
                                   WHERE deleted_at IS NOT NULL;

CREATE TABLE tag (
    id              TEXT PRIMARY KEY,
    name            TEXT NOT NULL,
    kind            TEXT NOT NULL DEFAULT 'user',
    CHECK (kind IN ('user','system')),
    UNIQUE (name, kind)
) STRICT;

CREATE TABLE asset_tag (
    asset_id        TEXT NOT NULL REFERENCES asset(id)  ON DELETE CASCADE,
    tag_id          TEXT NOT NULL REFERENCES tag(id)    ON DELETE CASCADE,
    PRIMARY KEY (asset_id, tag_id)
) STRICT, WITHOUT ROWID;

CREATE INDEX idx_asset_tag_tag ON asset_tag(tag_id, asset_id);
\end{lstlisting}

\subsection{Edit State}

\begin{lstlisting}[language=SQL, caption={Edit, history, and profile binding.}, label={lst:ddl-editfull}]
CREATE TABLE edit (
    asset_id         TEXT PRIMARY KEY REFERENCES asset(id) ON DELETE CASCADE,
    recipe_json      TEXT NOT NULL,
    recipe_version   INTEGER NOT NULL,
    recipe_hash      TEXT NOT NULL,
    source_recipe_id TEXT REFERENCES recipe(id) ON DELETE SET NULL,
    updated_at       REAL NOT NULL
) STRICT;

CREATE INDEX idx_edit_hash   ON edit(recipe_hash);
CREATE INDEX idx_edit_source ON edit(source_recipe_id)
                             WHERE source_recipe_id IS NOT NULL;

CREATE TABLE edit_history (
    id              INTEGER PRIMARY KEY AUTOINCREMENT,
    asset_id        TEXT NOT NULL REFERENCES asset(id) ON DELETE CASCADE,
    seq             INTEGER NOT NULL,
    recipe_json     TEXT NOT NULL,
    changed_key     TEXT,
    created_at      REAL NOT NULL
) STRICT;

CREATE UNIQUE INDEX idx_hist_asset_seq ON edit_history(asset_id, seq);

-- The (edit, profile version) binding. An edit references
-- every profile it was authored against, so a corrected profile can be
-- offered as an explicit opt-in rather than silently changing the image.
CREATE TABLE edit_profile_binding (
    asset_id        TEXT NOT NULL REFERENCES asset(id) ON DELETE CASCADE,
    profile_id      TEXT NOT NULL,
    profile_version INTEGER NOT NULL,
    role            TEXT NOT NULL,     -- 'negative'|'print'|'chemistry'
    PRIMARY KEY (asset_id, role),
    CHECK (role IN ('negative','print','chemistry')),
    FOREIGN KEY (profile_id, profile_version)
        REFERENCES stock_profile(id, profile_version) ON DELETE RESTRICT
) STRICT, WITHOUT ROWID;
\end{lstlisting}

The \texttt{ON DELETE RESTRICT} on the profile binding is deliberate and is the
enforcement point for the non-destructive promise of FR-12: a superseded
profile version cannot be deleted while any edit still references it, so
"clean up old profiles" cannot become "silently change the user's photographs."

\subsection{Recipes and Profiles}

\begin{lstlisting}[language=SQL, caption={Recipe and versioned stock profile.}, label={lst:ddl-recipefull}]
CREATE TABLE recipe (
    id              TEXT PRIMARY KEY,
    name            TEXT NOT NULL,
    recipe_json     TEXT NOT NULL,
    recipe_version  INTEGER NOT NULL,
    is_builtin      INTEGER NOT NULL DEFAULT 0,
    swatch_blob     BLOB,
    use_count       INTEGER NOT NULL DEFAULT 0,
    last_used_at    REAL,
    created_at      REAL NOT NULL,
    sort_order      INTEGER NOT NULL DEFAULT 0
) STRICT;

CREATE INDEX idx_recipe_used ON recipe(last_used_at DESC);

-- Primary key is (id, profile_version): superseded versions are retained,
-- not overwritten. Bundled profiles are re-inserted on launch only when the
-- (id, version) pair is absent.
CREATE TABLE stock_profile (
    id              TEXT NOT NULL,
    profile_version INTEGER NOT NULL,
    kind            TEXT NOT NULL,
    display_name    TEXT NOT NULL,
    params_json     TEXT NOT NULL,
    fit_status      TEXT NOT NULL DEFAULT 'E',
    is_builtin      INTEGER NOT NULL DEFAULT 1,
    superseded_by   INTEGER,
    installed_at    REAL NOT NULL,
    PRIMARY KEY (id, profile_version),
    CHECK (kind IN ('negative','print','mono','chemistry')),
    CHECK (fit_status IN ('M','P','E'))
) STRICT, WITHOUT ROWID;

CREATE INDEX idx_profile_current ON stock_profile(kind, id)
                                 WHERE superseded_by IS NULL;
\end{lstlisting}

\subsection{Rendition Cache and Export}

\begin{lstlisting}[language=SQL, caption={Rendition cache and export jobs.}, label={lst:ddl-rendfull}]
CREATE TABLE rendition (
    id              INTEGER PRIMARY KEY AUTOINCREMENT,
    asset_id        TEXT NOT NULL REFERENCES asset(id) ON DELETE CASCADE,
    cache_key       TEXT NOT NULL,
    size_class      TEXT NOT NULL,
    pixel_width     INTEGER NOT NULL,
    pixel_height    INTEGER NOT NULL,
    color_space     TEXT NOT NULL,
    file_relpath    TEXT NOT NULL,
    byte_size       INTEGER NOT NULL,
    created_at      REAL NOT NULL,
    last_access_at  REAL NOT NULL,
    CHECK (size_class IN ('thumb','grid','screen','full'))
) STRICT;

CREATE UNIQUE INDEX idx_rend_key ON rendition(cache_key);
CREATE INDEX idx_rend_lru        ON rendition(last_access_at);
CREATE INDEX idx_rend_asset      ON rendition(asset_id, size_class);

CREATE TABLE export_job (
    id              TEXT PRIMARY KEY,
    asset_id        TEXT NOT NULL REFERENCES asset(id) ON DELETE CASCADE,
    state           TEXT NOT NULL,
    long_edge       INTEGER,              -- NULL means native resolution
    format          TEXT NOT NULL,        -- 'heif'|'jpeg'|'tiff16'|'png'
    color_space     TEXT NOT NULL,        -- 'displayP3'|'srgb'|'rec2020pq'
    quality         REAL,
    include_sidecar INTEGER NOT NULL DEFAULT 0,
    destination     TEXT NOT NULL,        -- 'photos'|'files'|'share'
    error_domain    TEXT,
    error_code      INTEGER,
    enqueued_at     REAL NOT NULL,
    started_at      REAL,
    finished_at     REAL,
    CHECK (state IN ('queued','running','done','failed','cancelled')),
    CHECK (format  IN ('heif','jpeg','tiff16','png'))
) STRICT;

CREATE INDEX idx_export_queue ON export_job(state, enqueued_at)
                              WHERE state IN ('queued','running');
\end{lstlisting}

\subsection{Search}

Library search is full-text over the fields a photographer actually recalls:
roll name, note, tag, camera, and lens. It is an external-content FTS5 table so
the text is not duplicated, kept current by triggers rather than by application
code --- a search index maintained in application code is an index that goes
stale the first time a write path is added without remembering it.

\begin{lstlisting}[language=SQL, caption={FTS5 index and synchronization triggers.}, label={lst:ddl-fts}]
CREATE VIRTUAL TABLE asset_fts USING fts5(
    haystack,
    content = '',                 -- contentless; we store only the index
    tokenize = 'unicode61 remove_diacritics 2'
);

-- Rowid of asset_fts is the asset's integer surrogate, materialized here
-- because FTS5 requires an integer key and asset.id is a UUID string.
CREATE TABLE asset_fts_map (
    rowid           INTEGER PRIMARY KEY AUTOINCREMENT,
    asset_id        TEXT NOT NULL UNIQUE REFERENCES asset(id) ON DELETE CASCADE
) STRICT;

CREATE TRIGGER trg_asset_fts_ins AFTER INSERT ON asset BEGIN
    INSERT INTO asset_fts_map(asset_id) VALUES (new.id);
    INSERT INTO asset_fts(rowid, haystack)
    SELECT m.rowid,
           COALESCE(new.camera_model,'') || ' ' ||
           COALESCE(new.lens_model,'')   || ' ' ||
           COALESCE((SELECT name  FROM roll WHERE id = new.roll_id), '') || ' ' ||
           COALESCE((SELECT notes FROM roll WHERE id = new.roll_id), '')
    FROM asset_fts_map m WHERE m.asset_id = new.id;
END;

CREATE TRIGGER trg_asset_fts_del AFTER DELETE ON asset BEGIN
    INSERT INTO asset_fts(asset_fts, rowid, haystack)
    SELECT 'delete', m.rowid, '' FROM asset_fts_map m WHERE m.asset_id = old.id;
    DELETE FROM asset_fts_map WHERE asset_id = old.id;
END;
\end{lstlisting}

\subsection{Views}

Two views exist because the same three-way join appears in the library grid,
the filmstrip, and the export queue, and because a view is the cheapest place
to enforce that soft-deleted assets never appear in a list by accident.

\begin{lstlisting}[language=SQL, caption={Library views.}, label={lst:ddl-views}]
CREATE VIEW v_library AS
SELECT a.id, a.roll_id, r.name AS roll_name, a.captured_at, a.is_favorite,
       a.pixel_width, a.pixel_height, a.camera_model, a.lens_model,
       e.recipe_hash, e.recipe_version,
       (SELECT file_relpath FROM rendition
         WHERE asset_id = a.id AND size_class = 'grid'
         ORDER BY created_at DESC LIMIT 1) AS grid_relpath
FROM asset a
LEFT JOIN roll r ON r.id = a.roll_id
LEFT JOIN edit e ON e.asset_id = a.id
WHERE a.deleted_at IS NULL;

-- Cache pressure, read by the eviction task. Evicting thumb and grid last
-- is the stated policy, expressed where the query can see it rather than
-- in the application code that calls it.
CREATE VIEW v_evictable AS
SELECT rn.id, rn.file_relpath, rn.byte_size, rn.last_access_at,
       CASE rn.size_class WHEN 'full'   THEN 0
                          WHEN 'screen' THEN 1
                          WHEN 'grid'   THEN 2
                          ELSE 3 END AS evict_rank
FROM rendition rn
JOIN asset a ON a.id = rn.asset_id
ORDER BY evict_rank, rn.last_access_at;
\end{lstlisting}

\subsection{Integrity Constraints Not Expressible in DDL}

Three invariants are enforced in the repository layer because SQLite cannot
express them, and each is covered by a test that asserts the violation is
rejected:

\begin{enumerate}
\item \texttt{edit.recipe\_hash} must equal the canonical hash of
\texttt{edit.recipe\_json}. A trigger could recompute it only if SQLite could
canonicalize JSON, which it cannot; the repository computes it on write and a
consistency check verifies it during the periodic maintenance task.
\item \texttt{rendition.cache\_key} must incorporate the profile versions
bound in \texttt{edit\_profile\_binding}. Omitting them is the obvious
implementation and the wrong one: a corrected profile then serves stale
renditions from cache, and the user sees the old rendering until something
unrelated evicts it. The key is built in one function, and a test asserts that
bumping a profile version changes it.
\item \texttt{edit\_history.seq} must be contiguous from zero for each asset.
The unique index prevents duplicates but not gaps, which arise if a truncation
is interrupted; the maintenance task renumbers.
\end{enumerate}

\subsection{Expected Size}

For a library of 5\,000 assets with an average of 40 history entries each, the
database is approximately 95\,MB, of which 78\,MB is
\texttt{edit\_history.recipe\_json}. History is therefore the only table with a
retention policy: entries older than 90 days are pruned to the ten most recent
per asset by the maintenance task, which reduces the same library to about
22\,MB. Renditions are not counted here; they are files on disk, tracked by
\texttt{rendition.byte\_size} and capped at 3.5\,GB per Section~\ref{sec:schema}.
